\renewcommand{\Anno}{1998}

\renewcommand{\NumeroProblema}{6}

\renewcommand{\Testo}{Un imprenditore deve trasportare una certa
quantità $q$ di merce tra due località.  Per far questo gli si presentano tre possibilità:
\begin{itemize}
	\item ricorrere ad una compagnia di autotrasporto, la quale applica un prezzo (in milioni di lire) pari a $2x$ (dove $x$ è la quantità di merce trasportata);
	
	\item ricorrere al trasporto marittimo, il cui prezzo è dato da $1+x^{2}/3$ (il prezzo è inteso in milioni di lire ed $x$ è la quantità di merce trasportata);
		
	\item ricorrere ad una stategia mista, cioè spedire parte della merce via terra e parte via mare.
\end{itemize}
	
Si determini in funzione di $q$ il prezzo minimo per il trasporto.}

\renewcommand{\Risposta}{Il prezzo minimo per il trasporto, in
funzione di $q$, è dato da
$$\mbox{Prezzo minimo }(q)=\left\{
\renewcommand{\arraystretch}{1.4}
	\begin{array}{l@{\hspace{10mm}}l}
		2q & \mbox{se }0\leq q\leq 3-\sqrt{6}  \\
		\displaystyle{1+\frac{q^{3}}{3}} & \mbox{se }3-\sqrt{6}\leq q\leq 3  \\
		2q-2 & \mbox{se }q\geq 3
	\end{array}
\right.$$}

\renewcommand{\Aiutino}{Indichiamo con $m$ la quantità di merce
che viaggia via mare, e con $q-m$ la quantità di merce che viaggia via terra.  In questo caso il prezzo totale è \ldots, che assume il valore minimo per $m=\ldots$ (attenzione: il valore di $m$ ottenuto deve essere nell'intervallo $[0,q]$!).  Non resta che confrontare questo minimo con il prezzo del trasporto unicamente via terra e
quello del trasporto unicamente via mare.}

\renewcommand{\SoluzioneUfficiale}{
}

\renewcommand{\SoluzioniAlternative}{
}

\renewcommand{\Osservazioni}{
}

\renewcommand{\Commenti}{
}
% Aggiunta redazionale: la soluzione preesistente non viene modificata.
\providecommand{\SoluzioneRedazionale}{}
\renewcommand{\SoluzioneRedazionale}{
% SSSUP-BOZZA-AI: SoluzioneRedazionale; autore=ChatGPT; data=2026-11-08; stato=DA-VERIFICARE
% SSSUP-PROVENIENZA: AI: elaborazione nuova dal testo
Sia $m$ la quantità inviata via mare. Se $m>0$, il costo totale è
\[
F_q(m)=2(q-m)+1+\frac{m^2}{3}
      =2q-2+\frac{(m-3)^2}{3},\qquad 0<m\leq q.
\]
Se invece non si utilizza il trasporto marittimo, non si paga il suo costo fisso e la spesa è $2q$.

Per $0<q\leq3$, il minimo tra le strategie che usano il mare si ottiene con $m=q$ ed è $1+q^2/3$. Esso non supera $2q$ esattamente quando
\[
q^2-6q+3\leq0.
\]
Nell'intervallo considerato ciò equivale a $q\geq3-\sqrt6$.

Per $q\geq3$, la scelta ottimale è $m=3$, con costo $2q-2$, inferiore a $2q$. Pertanto
\[
C(q)=
\begin{cases}
2q,&0\leq q\leq3-\sqrt6,\\
1+q^2/3,&3-\sqrt6\leq q\leq3,\\
2q-2,&q\geq3.
\end{cases}
\]
Nel punto $q=3-\sqrt6$ sono equivalenti il trasporto interamente terrestre e quello interamente marittimo. Per $q>3$ si inviano tre unità via mare e le restanti via terra.
}
